\chapter{Tổng quan đề tài}
\label{chap:tong-quan}

\section{Lý do chọn đề tài}

Trong quá trình quản trị hệ thống mạng, kỹ thuật viên thường phải thao tác trực tiếp trên nhiều thiết bị như router, switch hoặc firewall. Việc cấu hình thủ công dễ mất thời gian, khó đồng bộ và có nguy cơ sai sót khi số lượng thiết bị tăng lên.

\section{Mục tiêu nghiên cứu}

Mục tiêu của đề tài là xây dựng phần mềm hỗ trợ quản lý và cấu hình thiết bị mạng sử dụng Python.

\section{Đối tượng và phạm vi nghiên cứu}

Đối tượng nghiên cứu gồm thiết bị mạng, giao thức kết nối quản trị, cơ sở dữ liệu lưu thông tin thiết bị và giao diện phần mềm quản lý.

\section{Phương pháp nghiên cứu}

Đề tài sử dụng phương pháp khảo sát yêu cầu, phân tích hệ thống, thiết kế cơ sở dữ liệu, xây dựng phần mềm và thử nghiệm trong môi trường mô phỏng như EVE-NG hoặc GNS3.

\section{Ý nghĩa khoa học và thực tiễn}

Kết quả đề tài giúp sinh viên tiếp cận bài toán tự động hóa mạng, kết hợp kiến thức mạng máy tính, lập trình Python và thiết kế phần mềm.
